\documentclass[11pt,letterpaper]{article}
\usepackage{cornell-notes}
\usepackage{needspace}

\title{Chapter 16: Network Application Protocols}
\author{}
\date{}

\setDocTitle{Chapter 16: Network Application Protocols}
\setDocSubtitle{Application Security Review}
\setCornellCollection{Software Security Assessment}
\setCornellUnitType{Chapter}
\setCornellUnitNumber{16}
\setCornellUnitTitle{Network Application Protocols}
\setCornellCompanionLabel{Study and audit reference}
\setCornellSourceRange{pp. 921--1005}
\begin{document}
\maketitle

\begin{center}
\small\color{Meta}\textbf{Coverage range:} pp. 921--1005 \quad\textbullet\quad \textbf{Format:} Cornell cue-and-notes
\end{center}
\vspace{0.25em}
\begin{CornellOverviewBox}{Review Orientation}
\textbf{Central problem.} Audit unfamiliar and familiar application protocols by reconstructing grammar, types, lengths, state, resource access, encoding rules, and trust assumptions from the bytes actually processed.
\textbf{Why it matters.} Security reviewers use this material to connect attacker-controlled conditions to violated invariants, consequential operations, and defensible remediation.
\textbf{Prerequisites.} Network byte order, parsing, integer safety, state machines, cryptographic protocol concepts, ASN.1, HTTP, and DNS basics.
\textbf{Assessment relationship.} It lifts low-level packet-validation techniques into structured application messages and complex standardized encodings.
\end{CornellOverviewBox}
\section{Learning Objectives}
\begin{itemize}
\item Explain the security significance of collect documentation.
\item Identify attacker influence and failed assumptions associated with identify unknown elements.
\item Trace security-relevant data or state through match data types.
\item Evaluate whether controls addressing data verification preserve the intended invariant.
\item Audit implementation and error paths related to system-resource access.
\item Distinguish safe and unsafe uses of header parsing.
\item Diagnose a failure involving resource access from concrete evidence.
\item Verify remediation and test coverage for utility functions and posted data.
\end{itemize}
\section{Cornell Cue-and-Notes}
\Needspace{8\baselineskip}
\subsection{General Protocol-Auditing Method}
\begin{longtable}{>{\raggedright\arraybackslash\columncolor{CornellCueBackground}}p{0.285\textwidth}|>{\raggedright\arraybackslash}p{0.625\textwidth}}
\rowcolor{Navy}\color{white}\textbf{Cue / Recall Prompt} & \color{white}\textbf{Technical Notes and Audit Reasoning} \\
\endfirsthead
\rowcolor{Navy}\color{white}\textbf{Cue / Recall Prompt} & \color{white}\textbf{Technical Notes and Audit Reasoning} \\
\endhead
\term{Collect documentation}\\[-0.1em]{\small What evidence defines intended protocol behavior?} & Specifications, implementation notes, packet captures, test suites, and source together reveal grammar and invariants. Documentation may omit extensions or differ from deployed behavior.\par\smallskip\textbf{Audit move:} Build a versioned protocol map and reconcile specifications with observed traffic and code. \\\hline
\term{Identify unknown elements}\\[-0.1em]{\small How can an undocumented protocol be reconstructed?} & Compare messages, vary one input at a time, locate length-like and discriminator fields, and trace decoder code. Separate observed facts from hypotheses until tests confirm them.\par\smallskip\textbf{Audit move:} Annotate captures with offsets, candidate types, dependencies, and confidence. \\\hline
\term{Match data types}\\[-0.1em]{\small Which language type receives each wire field?} & Wire widths, signedness, byte order, alignment, and repetition counts must map safely to host types. Conversions can create truncation or oversized allocation.\par\smallskip\textbf{Audit move:} Trace every field from raw bytes through conversion, validation, storage, and use. \\\hline
\term{Data verification}\\[-0.1em]{\small What relationships must a parser enforce?} & Validate outer bounds before inner fields, then check type-specific minimums, maximums, counts, lengths, terminators, and cross-field dependencies. Unknown values need explicit policy.\par\smallskip\textbf{Audit move:} Order validation so no read or arithmetic precedes the guarantee it requires. \\\hline
\term{System-resource access}\\[-0.1em]{\small How does protocol input reach files, commands, or state?} & Decoded fields may select paths, accounts, objects, algorithms, or backend requests. Parsing safety alone does not prevent authorization or injection flaws.\par\smallskip\textbf{Audit move:} Continue tracing validated fields into every privileged or persistent resource operation. \\\hline
\end{longtable}
\begin{CornellSummaryBox}{Section Summary}
Protocol review proceeds from documentation and byte structure to data types, validation, state, and resource use.
\end{CornellSummaryBox}
\Needspace{8\baselineskip}
\subsection{HTTP}
\begin{longtable}{>{\raggedright\arraybackslash\columncolor{CornellCueBackground}}p{0.285\textwidth}|>{\raggedright\arraybackslash}p{0.625\textwidth}}
\rowcolor{Navy}\color{white}\textbf{Cue / Recall Prompt} & \color{white}\textbf{Technical Notes and Audit Reasoning} \\
\endfirsthead
\rowcolor{Navy}\color{white}\textbf{Cue / Recall Prompt} & \color{white}\textbf{Technical Notes and Audit Reasoning} \\
\endhead
\term{Header parsing}\\[-0.1em]{\small How are field boundaries and duplicates interpreted?} & Line endings, whitespace, folding, delimiters, repeated fields, and length or transfer framing require one unambiguous policy. Intermediary disagreement can desynchronize message boundaries.\par\smallskip\textbf{Audit move:} Reject ambiguous framing and compare every hop's parsing rules. \\\hline
\term{Resource access}\\[-0.1em]{\small How does a request target map to an object?} & Decoding, normalization, virtual hosting, aliases, default resources, and filesystem rules transform a target into a resource. Authorization must apply after canonical mapping.\par\smallskip\textbf{Audit move:} Trace request target through all normalization and mapping steps. \\\hline
\term{Utility functions and posted data}\\[-0.1em]{\small Where do generic helpers hide assumptions?} & URL decoding, date parsing, token splitting, form decoding, and body handling are reusable attack surfaces. Body size and content type must be bounded before allocation and downstream parsing.\par\smallskip\textbf{Audit move:} Audit shared helpers once, then verify every caller supplies the promised preconditions. \\\hline
\end{longtable}
\begin{CornellSummaryBox}{Section Summary}
HTTP parsing is security-sensitive because flexible syntax, repeated fields, framing, and resource mapping can differ between intermediaries.
\end{CornellSummaryBox}
\Needspace{8\baselineskip}
\subsection{Key Management Protocol}
\begin{longtable}{>{\raggedright\arraybackslash\columncolor{CornellCueBackground}}p{0.285\textwidth}|>{\raggedright\arraybackslash}p{0.625\textwidth}}
\rowcolor{Navy}\color{white}\textbf{Cue / Recall Prompt} & \color{white}\textbf{Technical Notes and Audit Reasoning} \\
\endfirsthead
\rowcolor{Navy}\color{white}\textbf{Cue / Recall Prompt} & \color{white}\textbf{Technical Notes and Audit Reasoning} \\
\endhead
\term{Payload chains}\\[-0.1em]{\small How are nested payloads bounded?} & Payload type and length fields form chains whose termination, alignment, and total extent must remain inside the validated message. Unknown or repeated payloads need defined handling.\par\smallskip\textbf{Audit move:} Validate each link before advancing and enforce progress and aggregate limits. \\\hline
\term{Payload types}\\[-0.1em]{\small Can type confusion alter cryptographic meaning?} & Proposals, identities, nonces, keys, certificates, and notifications have distinct structures and state requirements. Accepting a valid structure in the wrong phase can weaken negotiation or authentication.\par\smallskip\textbf{Audit move:} Bind payload type and multiplicity to protocol state and negotiated algorithms. \\\hline
\term{Encryption vulnerabilities}\\[-0.1em]{\small Does encryption guarantee protocol security?} & Security also depends on authenticated identities, freshness, key derivation, algorithm negotiation, padding or integrity checks, and uniform failure handling. Parser or state flaws can bypass cryptographic intent.\par\smallskip\textbf{Audit move:} Trace exactly what is authenticated and bind it to session state and peer identity. \\\hline
\end{longtable}
\begin{CornellSummaryBox}{Section Summary}
Cryptographic protocols remain vulnerable when payload structure, negotiation, identity, and error behavior are parsed or bound incorrectly.
\end{CornellSummaryBox}
\Needspace{8\baselineskip}
\subsection{ASN.1 Encodings}
\begin{longtable}{>{\raggedright\arraybackslash\columncolor{CornellCueBackground}}p{0.285\textwidth}|>{\raggedright\arraybackslash}p{0.625\textwidth}}
\rowcolor{Navy}\color{white}\textbf{Cue / Recall Prompt} & \color{white}\textbf{Technical Notes and Audit Reasoning} \\
\endfirsthead
\rowcolor{Navy}\color{white}\textbf{Cue / Recall Prompt} & \color{white}\textbf{Technical Notes and Audit Reasoning} \\
\endhead
\term{Basic encoding rules}\\[-0.1em]{\small What makes tag-length-value parsing hazardous?} & Tags and lengths can be multi-byte, lengths may be indefinite in some rules, and constructed values nest recursively. Overflow, excessive depth, and length disagreement threaten memory and availability.\par\smallskip\textbf{Audit move:} Use bounded arithmetic, depth limits, and enclosing-slice validation at every node. \\\hline
\term{Canonical and distinguished rules}\\[-0.1em]{\small Why does canonical encoding matter?} & Canonical rules reduce multiple byte representations of the same value, which is crucial when signing, comparing, or caching encoded data. Accepting noncanonical forms can create identity or signature differentials.\par\smallskip\textbf{Audit move:} Canonicalize or strictly validate before security-sensitive comparison or signing. \\\hline
\term{Packed and XML encodings}\\[-0.1em]{\small How do alternate encoding rules change the attack surface?} & Bit-packed formats create intricate size arithmetic, while XML encodings add text parsing, entities, namespaces, and canonicalization. The abstract schema does not make every encoder equivalent.\par\smallskip\textbf{Audit move:} Audit each concrete encoding implementation and its resource limits independently. \\\hline
\end{longtable}
\begin{CornellSummaryBox}{Section Summary}
Structured encodings require safe handling of tag, length, nesting, canonical forms, and allocation before interpreting values.
\end{CornellSummaryBox}
\Needspace{8\baselineskip}
\subsection{Domain Name System}
\begin{longtable}{>{\raggedright\arraybackslash\columncolor{CornellCueBackground}}p{0.285\textwidth}|>{\raggedright\arraybackslash}p{0.625\textwidth}}
\rowcolor{Navy}\color{white}\textbf{Cue / Recall Prompt} & \color{white}\textbf{Technical Notes and Audit Reasoning} \\
\endfirsthead
\rowcolor{Navy}\color{white}\textbf{Cue / Recall Prompt} & \color{white}\textbf{Technical Notes and Audit Reasoning} \\
\endhead
\term{Names, servers, and zones}\\[-0.1em]{\small Which administrative boundary owns an answer?} & Resolvers query name servers across delegated zones and cache results. Authority, recursion, and bailiwick-like boundaries determine which data should influence future lookups.\par\smallskip\textbf{Audit move:} Model resolver roles, accepted sources, cache scope, and delegation transitions. \\\hline
\term{Message structure and records}\\[-0.1em]{\small How are counts and variable records validated?} & Header counts control repeated parsing, while each question and record has names, types, classes, lengths, and data. Aggregate counts and offsets must fit the packet and resource budget.\par\smallskip\textbf{Audit move:} Bound every count and record length against remaining bytes and total work. \\\hline
\term{Compressed names}\\[-0.1em]{\small Why can DNS name parsing loop or overrun?} & Compression pointers redirect parsing within the message and can form cycles, excessive chains, or invalid offsets. Output-name length also needs an independent bound.\par\smallskip\textbf{Audit move:} Track visited offsets, limit indirection, and validate expanded-name length. \\\hline
\term{Spoofing and cache integrity}\\[-0.1em]{\small What binds a response to a query?} & Source, identifiers, question match, timing, entropy, and protocol state influence acceptance. Weak matching or predictable state can let forged data enter caches.\par\smallskip\textbf{Audit move:} Require strict query-response matching and apply defenses appropriate to the trust model. \\\hline
\end{longtable}
\begin{CornellSummaryBox}{Section Summary}
DNS combines compressed names, counts, offsets, resource-record types, caching, and distributed trust.
\end{CornellSummaryBox}
\begin{CornellLegacyBox}{Historical Context}
Protocol standards and deployed extensions continue to evolve. Use the source's audit method while consulting current specifications for the exact protocol version under review.
\end{CornellLegacyBox}
\section{Security-Review Checklist}
\begin{CornellChecklist}
\item \textbf{Scoping}
Build a versioned protocol map and reconcile specifications with observed traffic and code.
\item \textbf{Attack-surface identification}
Annotate captures with offsets, candidate types, dependencies, and confidence.
\item \textbf{Input and data-flow analysis}
Trace every field from raw bytes through conversion, validation, storage, and use.
\item \textbf{Trust-boundary review}
Order validation so no read or arithmetic precedes the guarantee it requires.
\item \textbf{Candidate-point inspection}
Continue tracing validated fields into every privileged or persistent resource operation.
\item \textbf{Control-flow review}
Reject ambiguous framing and compare every hop's parsing rules.
\item \textbf{Resource and lifetime review}
Trace request target through all normalization and mapping steps.
\item \textbf{Error-path analysis}
Audit shared helpers once, then verify every caller supplies the promised preconditions.
\item \textbf{Mitigation validation}
Validate each link before advancing and enforce progress and aggregate limits.
\item \textbf{Evidence and reporting}
Bind payload type and multiplicity to protocol state and negotiated algorithms.
\item \textbf{Scoping}
Trace exactly what is authenticated and bind it to session state and peer identity.
\item \textbf{Attack-surface identification}
Use bounded arithmetic, depth limits, and enclosing-slice validation at every node.
\end{CornellChecklist}
\section{Chapter Synthesis}
Application-protocol review translates a wire format into a set of bounded slices, typed values, state transitions, and resource operations. The parser must reject ambiguity, prove progress, and preserve one interpretation across intermediaries. Prioritize recursive or compressed structures, cross-field length arithmetic, messages that allocate state, and values that select privileged resources or cryptographic identity.
\section{Self-Test Questions}
\begin{enumerate}
\item What evidence defines intended protocol behavior?
\item How can an undocumented protocol be reconstructed?
\item Which language type receives each wire field?
\item What relationships must a parser enforce?
\item How does protocol input reach files, commands, or state?
\item How are field boundaries and duplicates interpreted?
\item \textbf{Scenario:} A DNS decompressor follows pointers without recording visited offsets. What crafted input should be considered?
\item \textbf{Scenario:} Two HTTP intermediaries disagree about which framing header takes precedence. What class of vulnerability can result?
\end{enumerate}
\clearpage
\subsection*{Answer Key}
\begin{enumerate}
\item Specifications, implementation notes, packet captures, test suites, and source together reveal grammar and invariants. Documentation may omit extensions or differ from deployed behavior.
\item Compare messages, vary one input at a time, locate length-like and discriminator fields, and trace decoder code. Separate observed facts from hypotheses until tests confirm them.
\item Wire widths, signedness, byte order, alignment, and repetition counts must map safely to host types. Conversions can create truncation or oversized allocation.
\item Validate outer bounds before inner fields, then check type-specific minimums, maximums, counts, lengths, terminators, and cross-field dependencies. Unknown values need explicit policy.
\item Decoded fields may select paths, accounts, objects, algorithms, or backend requests. Parsing safety alone does not prevent authorization or injection flaws.
\item Line endings, whitespace, folding, delimiters, repeated fields, and length or transfer framing require one unambiguous policy. Intermediary disagreement can desynchronize message boundaries.
\item Pointers can form a cycle or excessive chain, causing nontermination or resource exhaustion. Validate offsets, track visits, cap indirection, and bound expanded output.
\item They may disagree about message boundaries, enabling request desynchronization or policy bypass. Enforce one unambiguous framing rule and reject conflicting inputs at every hop.
\end{enumerate}
\section{Key-Term Recap}
\begin{longtable}{>{\raggedright\arraybackslash}p{0.27\textwidth}p{0.65\textwidth}}
\toprule
\textbf{Term} & \textbf{Working meaning for review} \\
\midrule
\endfirsthead
\toprule
\textbf{Term} & \textbf{Working meaning for review} \\
\midrule
\endhead
Collect documentation & Specifications, implementation notes, packet captures, test suites, and source together reveal grammar and invariants. \\
Identify unknown elements & Compare messages, vary one input at a time, locate length-like and discriminator fields, and trace decoder code. \\
Match data types & Wire widths, signedness, byte order, alignment, and repetition counts must map safely to host types. \\
Data verification & Validate outer bounds before inner fields, then check type-specific minimums, maximums, counts, lengths, terminators, and cross-field dependencies. \\
System-resource access & Decoded fields may select paths, accounts, objects, algorithms, or backend requests. \\
Header parsing & Line endings, whitespace, folding, delimiters, repeated fields, and length or transfer framing require one unambiguous policy. \\
Resource access & Decoding, normalization, virtual hosting, aliases, default resources, and filesystem rules transform a target into a resource. \\
Utility functions and posted data & URL decoding, date parsing, token splitting, form decoding, and body handling are reusable attack surfaces. \\
Payload chains & Payload type and length fields form chains whose termination, alignment, and total extent must remain inside the validated message. \\
Payload types & Proposals, identities, nonces, keys, certificates, and notifications have distinct structures and state requirements. \\
\bottomrule
\end{longtable}
\section{One-Sentence Takeaway}
\begin{CornellExamBox}{One-Sentence Takeaway}
\textbf{A secure protocol parser turns hostile bytes into one bounded, canonical, state-valid interpretation before those values can allocate resources or select authority.}
\end{CornellExamBox}
\end{document}
